\documentclass[12pt]{article}
\usepackage[T2A]{fontenc}
\usepackage[utf8]{inputenc}
\usepackage[ukrainian]{babel}
\usepackage{hyperref}
\usepackage{geometry}
\geometry{a4paper, margin=1in}

\title{Software Requirements Specification (SRS)\\ \large Workout Volume \& Progressive Overload Tracker}
\author{System Analyst}
\date{\today}

\begin{document}

\maketitle
\tableofcontents
\newpage

\section{Вступ}

\subsection{Призначення системи}
Інформаційна система є вузькоспеціалізованим трекером силових тренувань, призначеним для автоматизації обліку механічної роботи (Volume Load) та алгоритмічного моніторингу прогресивного перевантаження (Progressive Overload). Основна категорія користувачів — атлети середнього та просунутого рівнів.

\subsection{Загальний опис}
Головна мета системи полягає у розв'язанні проблеми ручного підрахунку сумарного тоннажу та перерахунку одноповторного максимуму (1RM) у складних тренувальних програмах. Відсутність точної аналітики може призвести до стагнації або перетренованості. Ця система діє як інформаційний дашборд, який виконує всі математичні розрахунки та візуалізує Дельту прогресу, залишаючи експертні висновки самому атлету.

\subsection{Межі системи (Scope)}
\begin{itemize}
    \item Облік виключно силових тренувань (вправи, підходи, вага, повторення, суб'єктивне навантаження RPE 6.0--10.0).
    \item Розрахунок тренувального об'єму (Volume Load) для вправ з вільними вагами, тренажерами та власною вагою (включно з вправами з компенсацією ваги).
    \item Автоматичний "знімок" (snapshot) маси тіла користувача на момент проведення сесії для збереження історичної цілісності розрахунків.
    \item Гранулярне налаштування математичних моделей розрахунку 1RM на рівні кожної вправи (Epley як базова, Brzycki як альтернативна).
    \item Відображення інформаційних дашбордів із розрахунком Дельти (різниці Volume Load та 1RM) відносно останнього хронологічного виконання тієї ж вправи.
\end{itemize}

\subsection{Поза межами системи (Out of Scope)}
\begin{itemize}
    \item Трекінг кардіо-тренувань (біг, вело тощо).
    \item Ведення замірів пропорцій тіла та відсотка жиру (окрім загальної маси тіла).
    \item Трекінг харчування (калорії, макронутрієнти).
    \item Соціальні функції (шеринг, стрічка новин, підписки).
    \item Експертні системи: генерування автоматичних тригерів про "стагнацію" чи порад щодо зміни програми.
\end{itemize}

\section{Глосарій предметної області}
\begin{description}
    \item[Профіль користувача (User Profile)] Об'єкт у базі даних, який зберігає глобальні налаштування користувача, включно з його поточною масою тіла (оновлюється користувачем вручну).
    \item[Тренувальна сесія (Workout Session)] Логічний контейнер, що об'єднує набір виконаних вправ за один день. Під час створення сесії система автоматично фіксує поточну масу тіла користувача з його профілю (Body Weight Snapshot).
    \item[Вправа (Exercise)] Базова сутність силового тренування. Має один з типів: \textit{Standard} (Вільні ваги / Тренажери), \textit{Bodyweight} (Власна вага) або \textit{Assisted} (Із компенсацією ваги).
    \item[Підхід (Set)] Окрема спроба виконання вправи, що характеризується трьома основними метриками: вагою обтяження (Weight), кількістю повторень (Reps) та суб'єктивною оцінкою зусиль (RPE).
    \item[Розминочний підхід (Warm-up Set)] Підхід, який виконується для розігріву перед робочими підходами. Ігнорується системою при розрахунках сумарного Volume Load та максимального 1RM.
    \item[RPE (Rate of Perceived Exertion)] Шкала суб'єктивної оцінки важкості підходу від 6.0 до 10.0.
    \item[Volume Load (Тоннаж)] Кількісна метрика механічної роботи. Для стандартних вправ: \textit{Вага обтяження $\times$ Кількість повторень}. Для вправ із власною вагою: \textit{(Знімок маси тіла + Додаткова вага) $\times$ Кількість повторень}. Для вправ із компенсацією "Додаткова вага" записується як від'ємне значення.
    \item[1RM (Одноповторний максимум)] Розрахунковий показник максимальної ваги. Точні математичні моделі: Epley: $W \times (1 + \frac{R}{30})$; Brzycki: $W \times (\frac{36}{37 - R})$, де $W$ — вага, $R$ — кількість повторень.
    \item[Дельта (Delta)] Математична різниця показників (Volume Load або 1RM) поточного тренування порівняно з останнім хронологічним записом виконання цієї ж вправи.
\end{description}

\section{Вимоги до інтерфейсів та обмеження}
\subsection{Інтерфейси користувача}
Дизайн інтерфейсу розробляється за принципом Mobile-First. Система оптимізується для використання на смартфонах у портретній орієнтації під час тренування у спортивному залі.

\subsection{Проєктні обмеження (Constraints)}
Увесь життєвий цикл розробки суворо підпорядковується методології Specification-Driven Development (SDD). Документація вимог (цей SRS) є Єдиним джерелом істини (Single Source of Truth). Будь-які архітектурні чи логічні зміни повинні спочатку бути внесені до специфікації, і лише після цього реалізовуватися в кодовій базі.

\section{Специфічні вимоги}

\subsection{Функціональні вимоги}

\subsubsection*{Управління сесіями та підходами}
\begin{itemize}
    \item \textbf{REQ-F-001:} Система повинна надавати можливість створювати нову тренувальну сесію.
    \item \textbf{REQ-F-002:} У момент створення тренувальної сесії система повинна автоматично фіксувати поточну масу тіла користувача (snapshot), яка застосовуватиметься для всіх розрахунків у межах цієї сесії.
    \item \textbf{REQ-F-003:} Система повинна дозволяти додавати вправи до створеної сесії.
    \item \textbf{REQ-F-004:} Для кожної доданої вправи система повинна дозволяти фіксувати підходи (Sets) із зазначенням ваги, кількості повторень та RPE (у діапазоні 6.0--10.0).
    \item \textbf{REQ-F-005:} Система повинна надавати можливість позначити будь-який підхід як "Розминочний" (Is Warm-up).
\end{itemize}

\subsubsection*{Обчислення Volume Load}
\begin{itemize}
    \item \textbf{REQ-F-006:} Система повинна автоматично розраховувати показник Volume Load для кожного підходу. Сумарний Volume Load розраховується суворо для конкретного блоку доданої вправи (якщо вправу додано двічі в одну сесію — це два незалежні блоки, їхній Volume Load не плюсується).
    \item \textbf{REQ-F-007:} Під час розрахунку Volume Load система повинна повністю ігнорувати підходи, позначені як "Розминочні".
    \item \textbf{REQ-F-008:} Для стандартних вправ Volume Load підходу обчислюється як: \textit{Вага $\times$ Кількість повторень}.
    \item \textbf{REQ-F-009:} Для вправ із власною вагою Volume Load підходу обчислюється як: \textit{(Snapshot маси тіла + Додаткова вага) $\times$ Кількість повторень}.
    \item \textbf{REQ-F-010:} Система повинна дозволяти введення від'ємного значення у поле "Додаткова вага" та коректно зменшувати підсумковий Volume Load для вправ із компенсацією ваги.
\end{itemize}

\subsubsection*{Обчислення 1RM}
\begin{itemize}
    \item \textbf{REQ-F-011:} Система повинна автоматично обчислювати показник 1RM для кожного підходу, який не позначено як "Розминочний".
    \item \textbf{REQ-F-012:} Система повинна дозволяти користувачеві вибирати математичну модель для розрахунку 1RM (Epley або Brzycki) в налаштуваннях кожної конкретної вправи.
    \item \textbf{REQ-F-013:} Під час створення нової вправи система повинна автоматично встановлювати формулу Epley як алгоритм розрахунку 1RM за замовчуванням.
    \item \textbf{REQ-F-014:} Якщо кількість повторень у підході перевищує 10, система повинна розрахувати значення 1RM, але відобразити попередження (Soft Constraint) про знижену точність.
\end{itemize}

\subsubsection*{Аналітика та моніторинг}
\begin{itemize}
    \item \textbf{REQ-F-015:} Для кожної виконаної вправи система повинна обчислювати Дельту сумарного Volume Load відносно останнього хронологічного виконання цієї ж вправи.
    \item \textbf{REQ-F-016:} Система повинна обчислювати Дельту максимального 1RM відносно останнього хронологічного виконання цієї ж вправи.
    \item \textbf{REQ-F-017:} Система повинна надавати візуальний дашборд для кожної вправи, який відображає динаміку зміни Volume Load та максимального 1RM.
    \item \textbf{REQ-F-018:} Якщо система не знаходить попередніх записів для виконання вправи, Дельта не розраховується.
\end{itemize}

\subsection{Нефункціональні вимоги}
\begin{itemize}
    \item \textbf{REQ-NF-001:} Час виконання розрахунків та оновлення інтерфейсу для 1RM, Volume Load та Дельти не повинен перевищувати 100 мілісекунд.
    \item \textbf{REQ-NF-002:} Завантаження аналітичного дашборду не повинно перевищувати 1 секунду за умови підключення з пропускною здатністю мінімум 5 Мбіт/с та затримкою (ping) до 50 мс.
    \item \textbf{REQ-NF-003:} Система повинна здійснювати автоматичне локальне збереження (auto-save) даних підходу негайно після введення.
    \item \textbf{REQ-NF-004:} У разі критичної помилки або закриття застосунку дані активної сесії не повинні втрачатися і мають бути відновлені при наступному запуску.
    \item \textbf{REQ-NF-005:} Дизайн інтерфейсу повинен розроблятися за принципом Mobile-First (портретна орієнтація).
    \item \textbf{REQ-NF-006:} Інтерактивні елементи повинні мати мінімальний розмір зони натискання (Target Size) 44x44 dp згідно з вимогами Apple HIG.
    \item \textbf{REQ-NF-007:} Увесь життєвий цикл розробки повинен суворо підпорядковуватися методології SDD.
    \item \textbf{REQ-NF-008:} Документація вимог (SRS) є Єдиним джерелом істини (SSoT).
\end{itemize}

\section{Критерії прийняття та BDD-сценарії}

\subsection{Блок 1. Управління сесіями та підходами}
\textbf{Критерії прийняття:}
\begin{itemize}
    \item \textbf{AC1 (REQ-F-001, 002):} При створенні сесії записується дата, час та незмінний \texttt{body\_weight\_snapshot}.
    \item \textbf{AC2 (REQ-F-003, 004, 005):} Користувач може додати кілька вправ та підходів із зазначенням Ваги, Повторень, RPE та чекбоксу "Розминка".
\end{itemize}

\textbf{BDD-сценарій (BDD-001): Історична цілісність маси тіла}
\begin{verbatim}
Сценарій: Зміна поточної маси тіла не впливає на минулі тренування
Дано існуючу тренувальну сесію, створену місяць тому, зі снапшотом 80 кг
Коли користувач оновлює масу тіла в профілі до 85 кг
І система обчислює Volume Load для вправи з власною вагою з тієї сесії
Тоді для розрахунку має використовуватися значення снапшоту (80 кг)
І значення Volume Load цієї сесії залишається незмінним
\end{verbatim}

\subsection{Блок 2. Обчислення Volume Load}
\textbf{Критерії прийняття:}
\begin{itemize}
    \item \textbf{AC3 (REQ-F-006, 008):} Для стандартних вправ Volume Load дорівнює сумі (Вага $\times$ Повторення). Розрахунок відбувається строго в межах одного блоку вправи.
    \item \textbf{AC4 (REQ-F-007):} Підходи з прапорцем "Розминка" дають нульовий внесок у Volume Load.
    \item \textbf{AC5 (REQ-F-009, 010):} Для "Bodyweight" застосовується формула (Снапшот + Додаткова Вага) $\times$ Повторення. Допускається від'ємна вага.
\end{itemize}

\textbf{BDD-сценарій (BDD-002): Вправи з компенсацією (Гравітрон)}
\begin{verbatim}
Сценарій: Розрахунок об'єму для підтягувань із підтримкою
Дано поточну тренувальну сесію зі снапшотом маси тіла 80 кг
І вправу "Вправа з власною вагою"
Коли користувач записує підхід із "Додатковою вагою" -20 кг на 10 повторень
Тоді Volume Load має бути обчислений як (80 + (-20)) * 10
І система повинна зафіксувати результат 600 кг
\end{verbatim}

\textbf{BDD-сценарій (BDD-003): Ігнорування розминочних підходів}
\begin{verbatim}
Сценарій: Розминочні підходи виключаються з підрахунку тоннажу
Дано вправу "Присідання" в поточній сесії
Коли користувач виконує розминочний підхід з вагою 60 кг на 10 повторень
І виконує робочий підхід з вагою 100 кг на 5 повторень
Тоді сумарний Volume Load для вправи має становити 500 кг
\end{verbatim}

\subsection{Блок 3. Обчислення 1RM}
\textbf{Критерії прийняття:}
\begin{itemize}
    \item \textbf{AC6 (REQ-F-011, 012, 013):} 1RM автоматично розраховується для робочих підходів за обраною формулою (Epley за замовчуванням).
    \item \textbf{AC7 (REQ-F-014):} Якщо Повторення > 10, відображається індикатор "Soft Constraint".
\end{itemize}

\textbf{BDD-сценарій (BDD-006): Базовий розрахунок 1RM за формулою Epley}
\begin{verbatim}
Сценарій: Коректне застосування математичної моделі Epley
Дано робочий підхід із вагою (W) 100 кг та кількістю повторень (R) 5
І вправу налаштовано на використання формули Epley
Коли система обчислює показник 1RM
Тоді результат має дорівнювати 100 * (1 + 5/30) = 116.66 кг
\end{verbatim}

\textbf{BDD-сценарій (BDD-004): Soft Constraint при високій кількості повторень}
\begin{verbatim}
Сценарій: Інформування про можливу похибку розрахунку 1RM
Дано вправу з алгоритмом Epley
Коли користувач записує підхід з вагою 100 кг на 12 повторень
Тоді система розраховує 1RM за формулою Epley
І маркує результат попередженням про знижену точність (Soft Constraint)
\end{verbatim}

\subsection{Блок 4. Аналітика та моніторинг (Дельта)}
\textbf{Критерії прийняття:}
\begin{itemize}
    \item \textbf{AC8 (REQ-F-015, 016, 018):} Дельта є різницею між поточною та останньою хронологічною сесією тієї ж вправи.
    \item \textbf{AC9 (REQ-F-017):} Дашборд будує графік 1RM та Volume Load на основі історії.
\end{itemize}

\textbf{BDD-сценарій (BDD-005): Розрахунок Дельти}
\begin{verbatim}
Сценарій: Коректне порівняння прогресу з останньою сесією
Дано історію тренувань: "Станова тяга" 1 Вересня (Volume Load: 3000 кг)
І "Станова тяга" 8 Вересня (Volume Load: 3200 кг)
Коли користувач завершує сесію 15 Вересня із Volume Load 3300 кг
Тоді Дельта розраховується лише на основі сесії від 8 Вересня
І значення Дельти має становити +100 кг
\end{verbatim}

\textbf{BDD-сценарій (BDD-007): Рендеринг графіка на дашборді}
\begin{verbatim}
Сценарій: Дашборд успішно будує графік на основі історичних даних
Дано вправу, яка має щонайменше 2 збережені тренувальні сесії в історії
Коли користувач відкриває аналітичний дашборд цієї вправи
Тоді система повинна успішно відрендерити лінійний графік
І на графіку повинні бути точки, що відповідають значенням 1RM та Volume Load з цих сесій
\end{verbatim}

\end{document}
